\chapter{Experimental Methodology}
\label{ch:experiment}

This chapter defines the staged empirical protocol used after the Vericult architecture is frozen. The design separates development and stress testing, backbone selection, realistic confirmatory evaluation and external validation. This separation is necessary because the project has already used the original PLT set extensively during engineering; those observations should not be treated as an untouched final test set.

\section{Evaluation Strategy}

The empirical program has three roles:
\begin{enumerate}
    \item \textbf{Development and stress testing}: use the existing PLT v1 material to expose architectural weaknesses, validate failure handling and select the final locally runnable semantic backbone.
    \item \textbf{Realistic confirmatory evaluation}: construct PLT v2 from the same 30 prompts but with responses generated by stronger contemporary models, freeze those candidates before evaluation, and assess Vericult without further semantic tuning.
    \item \textbf{External validation}: evaluate the frozen verifier on deterministic samples from published cultural-evaluation resources, preserving each source's native task structure where possible.
\end{enumerate}

This design supports both the controlled Best-of-4 question and the intended single-response verification question. Best-of-4 tests relative preference selection; candidate-level labels and scores test the primary question: \emph{is this answer culturally appropriate, partially appropriate or culturally inappropriate?}

\section{PLT v1: Development and Stress-Test Set}

The project contains 30 prompts, \texttt{PLT001}--\texttt{PLT030}, with four responses per prompt. The original candidates were generated during the red-team and development phase and contain a comparatively high prevalence of obvious or exaggerated cultural problems. They are therefore useful for testing whether the verifier notices material cultural failures, but they are not treated as representative of the quality produced by current stronger aligned models.

The v1 set remains scientifically useful as development data. It was the dataset on which the initial full LIVE attempt exposed target-extraction reliability problems. Human annotation for this set is already frozen and remains useful for descriptive comparison and backbone selection. However, final claims about realistic model behavior are not based on v1 alone.

\section{Frozen Human Reference for v1}

Five independent annotators evaluated all 30 prompt sets. Contrary to an earlier planned protocol, the realized survey did not offer \emph{No clear winner} or a 1--5 confidence scale. Instead, each candidate response received an absolute appropriateness rating on a three-point scale, followed by a forced A/B/C/D choice for the best response. This realized protocol is authoritative for analysis.

The frozen export contains 600 candidate-level ratings (30 prompts $\times$ 4 candidates $\times$ 5 annotators) and 150 forced winner votes. Across the 30 prompts, 17 achieved a winner majority of at least three of five annotators, while 13 remained ambiguous under that majority rule. Ten prompts had ties for the highest vote count and one prompt was unanimous. Fleiss' $\kappa$ for winner votes is 0.090494, indicating low inter-annotator agreement and reinforcing that the aggregate is a human reference rather than objective cultural ground truth. The candidate-rating distribution is 202 ratings of 1, 207 ratings of 2 and 191 ratings of 3.

The forced-winner design is especially important when interpreting cases such as PLT020, where all candidates can be rated poorly but one still receives the most winner votes. Absolute candidate appropriateness and relative Best-of-4 preference are therefore analyzed separately.

\section{PLT v2: Realistic Strong-Model Benchmark}

PLT v2 reuses the exact 30 PLT prompts but replaces all four candidate responses. Its purpose is to test Vericult when candidate answers are fluent and plausibly useful rather than intentionally failure-heavy. Candidate generation is performed before Vericult is run on v2 and without access to verifier scores, human labels or D01--D10 winner signals.

The generator set is chosen to balance contemporary answer quality, model-family diversity and practical cost. The final generation manifest records exact model identifiers, provider or local runtime, generation date, temperature, output-token limit and raw outputs. Each prompt receives one response from each frozen generator; source models are deterministically shuffled into A--D using a fixed seed. No response is selected because it appears culturally better or worse. Only technical quality gates such as nonempty successful generation are permitted.

The candidate generators are not verifier components. A frontier API model may be used for one generator while other candidates are produced locally; this is treated as benchmark construction rather than a dependency of Vericult. The exact v2 generator set is frozen before any v2 Vericult prediction is produced and is reported in the final generation manifest.

\section{Backbone Selection Before Final Freeze}

The Vericult architecture is fixed before backbone selection. Candidate backbones are selected from research-relevant open instruction-tuned models that can run on the available local Apple-silicon workspace. The selection criterion is not raw benchmark popularity alone. Each candidate backbone is evaluated on a small v1 diagnostic subset chosen before model outcomes are inspected, using the same architecture and retrieval configuration.

Backbone selection considers four quantities: (1) structural and technical completion rate, (2) agreement with the already frozen v1 human reference where comparison is available, (3) stability of dimension and target outputs under temperature zero, and (4) wall-clock feasibility for the final experiment. The chosen backbone and exact Ollama model digest are then frozen. v2 and external-validation results are not used for backbone selection.

CARB provides empirical motivation for including Qwen-family models in this candidate pool because its experiments report strong cultural reward-model and generative-evaluator performance from Qwen-based systems \parencite{zhang2026carb}. The exact locally runnable variant is nevertheless selected empirically rather than assumed to inherit the performance of a much larger model.

\section{Final Vericult Freeze}

After architecture validation and backbone selection, the following are frozen together: code commit, pipeline version, prompt templates, D01--D10 rubric, model identifier and digest, temperature, retry policy, target budget, retrieval depth, maximum rounds, provider configuration, leakage exclusions, decision rules and statistical-analysis seed. No semantic tuning is permitted after the first v2 prediction is produced.

The intended default constants remain temperature 0, at most three material targets, two initial neutral verification questions per retrievable target, \texttt{top\_k}=3 accepted documents per query, and at most two retrieval rounds. Any final change to these values occurs before v2 execution and is recorded in the final manifest.

\begin{table}[htbp]
\centering
\caption{Elements frozen before confirmatory evaluation.}
\begin{tabular}{p{5.0cm}p{8.5cm}}
\toprule
\textbf{Component} & \textbf{Frozen information} \\
\midrule
Implementation & Git commit, package and pipeline versions \\
Semantic model & Model ID, local digest, provider, temperature and retry policy \\
Cultural framework & D01--D10 rubric, anchors and aggregation rule \\
Semantic contracts & Prompt-template hashes and structured output schemas \\
Retrieval & Provider, \texttt{top\_k}, search depth, maximum rounds and exclusions \\
Evaluation & v2 candidate manifest, permutation seeds, baseline versions and statistical plan \\
\bottomrule
\end{tabular}
\end{table}

\section{LIVE Acquisition and REPLAY}

For evaluation sets requiring web evidence, the frozen configuration first runs in LIVE mode to create evidence schedules and retrieval snapshots. Those artifacts are hashed and archived. REPLAY then reuses the same evidence without live retrieval and serves as the primary Vericult result where replay is technically appropriate. A missing or inconsistent schedule is reported as a failure rather than silently replaced by new search.

Because v2 and external data are evaluated only after the semantic freeze, LIVE evidence can change with the public web without creating an opportunity for post-result prompt tuning. Retrieval dates and source hashes remain part of the reproducibility archive.

\section{Comparison Systems}

\subsection{Skywork Reward Model}

The primary reward-model baseline is \texttt{Skywork/Skywork-Reward-V2-Qwen3-4B}, revision \texttt{fd958fe} \parencite{liu2026skyworkrewardv2}. It receives only the prompt and one candidate response and remains independent of Vericult, D01--D10, retrieval and human labels. Four raw rewards induce a Best-of-4 preference. An exact maximum-score tie maps to \texttt{no\_clear\_winner}; raw reward magnitudes are not interpreted as calibrated cultural probabilities.

\subsection{Direct LLM Judge}

A direct judge receives the prompt and four candidate responses in a deterministic permutation, uses the D01--D10 rubric, performs no retrieval, and selects A/B/C/D or a selective outcome. Whenever feasible, the direct judge uses the same frozen semantic backbone as Vericult so that the comparison tests architectural structure rather than merely comparing a small verifier model with a stronger judge.

\begin{table}[htbp]
\centering
\caption{Final comparison systems and their intended roles.}
\begin{tabular}{p{3.4cm}p{4.4cm}p{5.8cm}}
\toprule
\textbf{System} & \textbf{Primary signal} & \textbf{Role} \\
\midrule
Vericult & D01--D10, retrieval, frozen evidence and contextual fallback & Proposed transparent verifier \\
Skywork Reward V2 & Learned scalar preference reward & Independent general RM baseline \\
Direct LLM judge & One rubric-guided no-retrieval judgment & Tests whether full architecture adds value over direct judging \\
Human reference & Candidate ratings and forced winner votes & Operational comparison reference, not universal cultural ground truth \\
\bottomrule
\end{tabular}
\end{table}

\section{CARB and Other External Validation Sources}

External evaluation is not limited to the project's own 30 prompts. After the final Vericult freeze, deterministic source-backed samples are drawn from complementary published research. CARB is used in its native Best-of-N preference setting when official author-released item-level data are available. This permits direct qualitative and, where compatible, quantitative comparison with the cultural reward-model findings reported by CARB \parencite{zhang2026carb}.

The preregistered external protocol additionally includes SafeWorld for culturally and legally situated safety, CultureLLM and World Values Survey material for plural values, and ScopeBench for universal-versus-specific cultural scope. Recent free-text resources such as FrameNet-Cultures and ExCAM are additionally relevant for validating arbitrary response assessment because they move beyond closed-form cultural QA \parencite{yin2024safeworld,li2024culturellm,jamshidi2026framenet,leiter2026excam}. External items are never used to tune the final frozen system.

Where a source does not natively provide four candidate answers, it is evaluated at candidate level rather than artificially converted into Best-of-4. This preserves the distinction between relative preference ranking and Vericult's primary single-response classification task.

\section{Primary and Secondary Metrics}

Two families of metrics are reported.

\paragraph{Absolute candidate assessment.}
Where candidate-level human or benchmark labels are available, Vericult's three-way cultural label and 0--100 score are compared with the reference. Ordinal association, class agreement and false endorsement or false rejection are reported where meaningful. Evidence coverage and contextual-fallback usage are diagnostics, not substitutes for correctness.

\paragraph{Best-of-4 preference assessment.}
For items with a resolved human or benchmark winner, exact winner agreement is reported. Exact ties, no-acceptable-candidate outcomes and technical failures are reported separately. For v1, the forced human winner is not interpreted as proof that the selected candidate is absolutely acceptable; candidate-level ratings are used to expose such cases.

Additional engineering metrics include completion rate, failures by pipeline stage, retrieval calls, evidence coverage, fallback rate and wall-clock runtime. These metrics are important because a verifier that is accurate only when it completes is not practically equivalent to a reliable evaluation layer.

\section{Statistical Analysis}

Uncertainty is analyzed at item level because the same prompts are evaluated by all systems. Bootstrap confidence intervals resample whole prompts with replacement \parencite{efron1993bootstrap}. Paired binary correctness comparisons use exact McNemar tests when discordant counts are small \parencite{mcnemar1947correlated}. Effect sizes and discordant counts accompany $p$-values.

Human agreement is reported rather than hidden. Fleiss' $\kappa$ is available for the complete five-rater v1 winner matrix \parencite{fleiss1971kappa}, while candidate-level three-point ratings can additionally be analyzed with an ordinal agreement statistic if used in the final results chapter. Small per-dimension or per-source cells are treated descriptively rather than overinterpreted.

\section{Execution Order and Reproducibility}

The final sequence is fixed:
\begin{enumerate}
    \item finish architecture CI and development diagnostics;
    \item select and freeze the verifier backbone on v1 development material;
    \item generate and freeze v2 candidates;
    \item freeze the v2 human reference before inspecting machine outcomes;
    \item acquire LIVE evidence;
    \item run primary REPLAY where applicable;
    \item run Skywork and direct-judge baselines;
    \item evaluate frozen external samples;
    \item compute pre-declared metrics and paired statistics; and
    \item only then select trace-backed cases for qualitative discussion.
\end{enumerate}

The separation between development and confirmatory phases is more important than preserving an earlier nominal final manifest that was created before reliability defects were discovered. The thesis reports the actual final frozen commit and manifest used for confirmatory evaluation, while earlier manifests and failed attempts remain engineering history rather than primary results.
